\chapter{Skiljet}

\begin{motto}
Den dagen ein mann teikna det ein annan skulle lage,\\ var designaren fødd, og avstanden til følgja med han.
\end{motto}

Det som skil designfaget frå handverket er ikkje maskina. Maskinar fanst lenge før design, og handverkarar har alltid brukt verktøy. Det som skil dei er ei deling: delinga mellom den som tenkjer ut forma og den som lagar tingen. I handverket budde desse to i éin kropp; den industrielle revolusjonen skilde dei, sette dei i kvar sin del av prosessen, ofte i kvar sin del av bygningen, og fødde dermed ein ny skikkelse som ikkje hadde funnest før: mannen som teiknar det andre skal lage. Det er denne skikkelsen, formgjevaren, designaren, historia vår eigentleg handlar om, og det fyrste vi må sjå er at han vart til som ein \emph{funksjon i ei arbeidsdeling}, ikkje som ein berar av eit nytt fag. Han fanst før det fanst noko fag å gje honom. Tomrommet under faget vart grave ut i same rørsla som skapte behovet for det.

Adam Smith hadde alt skildra prinsippet før det fekk konsekvensar for forma. I knappenålsfabrikken hans var arbeidet delt opp: ein mann drog tråden, ein annan retta han, ein tredje kappa, ein fjerde spissa. Kvar arbeidar gjorde éi rørsle om att og om att, og fabrikken som heilskap produserte meir enn dei same mennene kunne gjort kvar for seg. Arbeidsdelinga var produktiv; det er ikkje til å diskutere. Men legg merke til kva ho gjorde med kunnskapen. Mannen som spissa nåla heile dagen, visste korleis ein spissar ei nål, og ingenting meir. Han kjende ikkje tingen han var med på å lage; han kjende sin del av han. Den heile kunnskapen om nåla, frå tråd til ferdig produkt, fanst no ikkje lenger i nokon einaste arbeidar. Han var spreidd ut over mange, og samla att berre i éin: han som hadde planlagt delinga. Tenkinga om heilskapen flytta oppover og ut, vekk frå hendene som gjorde arbeidet.

\section{Frå verkstad til fabrikk}

Sjå konkret på kva som skjedde med ein gjenstand i overgangen frå verkstad til fabrikk, for det abstrakte prinsippet får fyrst tyngd når ein følgjer ein ting gjennom det. I verkstaden vart ein stol til under hendene til ein snikkar som heldt heile prosessen: han valde treet, han forma kvar del, han sette saman, og avgjerdene om form vart tekne undervegs, i materialet, av den same dømekrafta som førte verktøyet. Forma og produksjonen var ikkje to fasar; dei var same rørsle. I fabrikken vart dette teke frå kvarandre. Forma vart bestemt på førehand, i ei teikning, ein modell, ein mal, og produksjonen vart sett til å realisere ei form som alt var avgjort utanfor verkstaden. Arbeidaren ved maskina tok ingen formavgjerder. Han utførte ei form ein annan hadde fastsett.

Adrian Forty har vist, klårare enn nokon, at denne skilnaden ikkje fyrst og fremst var teknisk, men sosial og økonomisk: forma vart skild frå produksjonen fordi det tente kapitalen å gjere det \autocite{forty1986}. Den som eigde fabrikken, ville ikkje vere prisgjeven den tause kunnskapen til ein dyr handverkar som kunne gå sin veg og ta kunnskapen med seg. Ved å flytte forma ut av handa og inn i teikninga, gjorde eigaren kunnskapen om forma til noko han kunne eige, kontrollere og skilje frå arbeidaren. John Heskett har dokumentert korleis dette skapte ein heilt ny type arbeid, teiknaren, mønsterlagaren, modelløren, folk som vart tilsette nettopp for å produsere forma som så vart send vidare til verkstadgolvet for utføring \autocite{heskett1980}. Dette er designaren i embryo. Han er fødd, ikkje som ein kunstnar med ein visjon, men som ein funksjonær i ei arbeidsdeling: han som leverer forma slik at andre kan lage tingen utan å kunne forma sjølv.

David Pye sitt skilje mellom risikoens og visseheitens handverk får her sin historiske dato. Fabrikken er visseheitens handverk sett i system: resultatet er sikra på førehand, av malen, av maskina, av teikninga, slik at den som utfører ikkje lenger treng dømekraft. Og det var heile poenget. Den industrielle revolusjonen var, sett frå kunnskapssida, eit stort prosjekt for å fjerne risiko frå produksjonen ved å fjerne dømekrafta frå arbeidaren, og dømekrafta var det handverket sin kunnskap budde i. Det som vart att hjå arbeidaren, var rørsla utan kunnskapen. Det som flytta til teiknaren, var kunnskapen utan rørsla. Mennesket der dei to ein gong var eitt, fanst ikkje lenger i produksjonen.

\section{Teiknaren får eit yrke, men ikkje eit fag}

Sjølve skikkelsen som steig fram av skiljet, er hovudpersonen i resten av boka, og opphavet hans avgjer kva han kunne bli. Den industrielle teiknaren var ikkje ein arvtakar av handverksmeisteren; han var noko nytt, og det nye var presis at han hadde forma utan å ha lagingane. I tekstilindustrien voks det fram mønsterteiknarar som leverte design til vevstolar dei aldri rørte; i keramikken modellørar som forma originalen som så vart kopiert i tusental; i metallvareindustrien teiknarar som leverte snittet som maskina skulle realisere. Penny Sparke har peika på at dette er det yrkesmessige opphavet til det vi i dag kallar designaren: ikkje kunstnaren som steig ned i industrien, men funksjonæren som industrien sjølv skapte for å skilje formavgjerda frå produksjonen \autocite{sparke2004}. Yrket kom før kunsten, og lenge før faget.

Og her ligg eit avgjerande poeng for heile soga. Teiknaren fekk eit \emph{yrke}, ein plass i lønningslista, ein funksjon i kjeda, etter kvart ein tittel, men han fekk ikkje med det eit \emph{fag} i den meininga eit fag ber: ein delt kunnskap om kva han dreiv med, ein teori om forma han produserte, ein standard han kunne målast mot som ikkje var arbeidsgjevaren sin smak eller marknaden sin dom. Jonathan Woodham har vist korleis designaren gjennom det meste av industrialderen var ein anonym arbeidstakar, definert av produksjonen han tente og ikkje av nokon eigen kunnskap, og korleis sjølve ideen om designaren som ein autonom fagperson med eit eige kompetansegrunnlag er ei sein og dels etterkonstruert forteljing \autocite{woodham1997}. Yrket fanst i hundre år før faget freista å gje det ein grunn. Når vi seinare ser rørsle etter rørsle prøve å heve designaren frå funksjonær til fagperson, må vi hugse at dei arbeidde med eit materiale som frå fyrste stund var definert ovanfrå, av kva industrien trong, og ikkje nedanfrå, av ein kunnskap som var teiknaren sin eigen. Eit yrke utan eit fag er nett det skiljet etterlét: ein som gjer formarbeid kompetent, utan å eige forma han gjer.

\section{Wedgwood, og delinga sett i drift}

Det abstrakte prinsippet får namn, stad og dato i Staffordshire. Då Josiah Wedgwood opna fabrikken sin Etruria i 1769, sette han i system det knappenålsfabrikken berre hadde skissert: han delte leirpottemakeriet, eit av dei eldste samla handverka som finst, opp i ei rekkje åtskilde operasjonar utførte av åtskilde hender. Der pottemakaren før hadde halde heile gjenstanden frå klump til glasur, vart arbeidet no spreidd på dreiarar, formarar, handtak-festarar, målarar, glasurarbeidarar og omnsarbeidarar, kvar med si avgrensa rørsle. Wedgwood skreiv det sjølv, i eit brev historikarane stadig vender attende til: målet var å «make such machines of the men as cannot err». Å gjere slike maskiner av mennene at dei ikkje kan feile. Setninga er heile skiljet i ni ord. Feilen, risikoen i Pye si meining, skulle fjernast frå arbeidaren ved å fjerne dømekrafta frå han, og med dømekrafta fylgde den einaste formkunnskapen handverket nokon gong hadde ått ut av kroppen hans.

Og kvar flytta kunnskapen? Til teiknerommet og modellrommet. Wedgwood kjøpte forma utanfrå, frå folk som aldri rørte ein omn: han engasjerte biletmakaren John Flaxman til å levere relieff og motiv som så vart kopierte i serie av hender som ikkje hadde teikna dei. Her ser ein figuren reint. Flaxman teikna; arbeidarane ved Etruria realiserte; og mellom dei to gjekk det ei deling som ikkje fanst i pottemakarverkstaden. Forma var no ei vare Wedgwood kunne kjøpe, eige og setje i produksjon uavhengig av den tause kunnskapen til nokon einskild handverkar, som var heile vitsen, fordi ein handverkar med kunnskapen i handa kunne gå sin veg og ta han med seg, medan ei teikning blir liggjande. Adrian Forty har lese Wedgwood som det tidlege mønsterdømet på at forma vart skild frå produksjonen fordi det tente eigaren, ikkje fordi det tente gjenstanden \autocite{forty1986}. John Heskett peikar på det same: at det er her, i den keramiske industrien i siste fjerdedel av 1700-talet, at ein av dei fyrste klart utskilde «designar»-funksjonane i moderne meining trer fram, ikkje som kunstnar, men som leverandør av form til ein produksjon som hadde slutta å kunne forme sjølv \autocite{heskett1980}.

\begin{kasus}{«Maskiner av menn» – Etruria 1769}
Wedgwood sitt eige ordval er meir avslørande enn noka seinare analyse. Han ville lage «such machines of the men as cannot err», og han lukkast så godt at namnet hans framleis står på porselen i dag, to og eit halvt hundreår seinare. Men sjå kva prisen var, og kven som betalte han. Pottemakaren som før heldt heile kunnskapen, vart redusert til ein som festa handtak heile dagen og ikkje visste, ikkje trong vite, ikkje fekk vite, korkje korleis leira vart blanda eller korleis omnen vart styrt. Den heile kompetansen som hadde butt i éin kropp, vart pulverisert og spreidd, og kunne ikkje lenger samlast att av nokon på golvet. Wedgwood fjerna feilen frå produksjonen ved å fjerne kunnskapen frå arbeidaren, og flytta han til teiknerommet der Flaxman sat. Det er den industrielle forma sin fødsel som administrativt faktum: ikkje ein ny estetikk, men ei ny eigedomsordning over kunnskapen om form. Den som vil forstå kvifor designaren frå fyrste stund var ein funksjonær og ikkje ein fagperson, treng ikkje sjå lenger enn til kven som teikna og kven som festa handtak ved Etruria.
\end{kasus}

Det avgjerande er at denne delinga ikkje var ein dekadanse, men ein suksess, og det er nett det som gjer henne uimotståeleg og varig. Etruria produserte vakrare og jamnare gods enn det gamle verkstaden kunne, til ein pris fleire kunne betale, og selde det over heile Europa. Delinga \emph{verka}. Penny Sparke har understreka at heile den seinare designhistoria må lese denne suksessen ærleg: forma vart ikkje verre av å bli skild frå handa, ho vart på mange vis betre kontrollert og meir tilgjengeleg \autocite{sparke2004}. Men ho vart heimlaus. Kunnskapen om kvifor forma var god, budde no ikkje lenger i nokon, korkje i arbeidaren, som berre hadde si rørsle, eller i Flaxman, som berre hadde teikninga si, eller i Wedgwood, som hadde rekneskapen. Forma var blitt eit produkt av eit system, og ingen ledd i systemet eigde kunnskapen om henne som heilskap. Det er denne heimløysa, ikkje noko estetisk forfall, som er det verkelege tomrommet faget seinare vart kalla inn for å fylle.

\section{Form skild frå produksjon, og frå følgje}

To skiljer skjer her samstundes, og det er viktig å halde dei frå kvarandre, fordi det andre er det boka kjem til å forfølgje resten av vegen. Det fyrste er skiljet mellom form og produksjon: formgjevaren teiknar, andre lagar. Det andre, og alvorlegare, er skiljet mellom form og \emph{følgje}. Så lenge handverkaren laga heile tingen for ein kjend kunde i ein kjend stad, stod han andsynes konsekvensane av arbeidet sitt. Laga han ein dårleg stol, fekk han det att frå grannen som kjøpte han. Brukte han dårleg tre, song det i hans eige rykte. Forma, produksjonen og følgja budde i same nærleik, i same krins av menneske som kjende kvarandre. Handverkaren bar konsekvensen av forma si fordi han var der konsekvensen fall.

Den industrielle forma fell ein annan stad. Teiknaren i Birmingham teiknar ei vare som vert laga av hender han aldri ser, seld i ein marknad han aldri går i, brukt av menneske han aldri møter, og kasta ein stad han aldri kjem. Mellom forma og følgja ligg det no fabrikkar, grossistar, butikkar, eit halvt kontinent. Formgjevaren ser ikkje lenger kva forma hans gjer i verda, fordi han er strukturelt skild frå staden der ho gjer det. Dette er ikkje ein moralsk svikt hjå den enkelte teiknaren; det er ein eigenskap ved sjølve arbeidsdelinga. Ho plasserer den som bestemmer forma så langt frå den som ber konsekvensen at konsekvensen ikkje lenger kan informere forma. Frank Trentmann har kartlagt korleis denne avstanden voks i takt med at varene vart fleire, billegare og meir flyktige: ein heil sivilisasjon av ting laga av folk som aldri ville sjå kva dei vart til \autocite{trentmann2016}.

Det er her, ikkje seinare, ikkje i overvakingskapitalismen, men her, ved industrialiseringa sjølv, at den figuren oppstår som boka kallar at \emph{skaden ber andre}. Formgjevaren haustar forma; følgja av forma fell på arbeidaren som slit ved maskina, på kjøparen som får ei vare laga for å seljast og ikkje for å vare, på staden der ho vert kasta. Og fordi formgjevaren ikkje ber følgja sjølv, treng ho ikkje gå inn i avgjerda hans. Han kan optimere forma langs den eine aksen oppdragsgjevaren bryr seg om, pris, sal, framtoning, utan at dei andre aksane melder seg, fordi dei melder seg ein annan stad. Eit fag med eit fundament ville hatt eit språk for desse andre aksane, ein måte å tvinge følgja attende inn i forma. Designfaget kjem aldri til å byggje eit slikt språk, og grunnen er at faget vart fødd nett i det øyeblikket avstanden vart oppretta. Det arva ikkje berre handverkets tause kunnskap; det arva ein arbeidsdeling som plasserte formgjevaren utanfor rekkjevidde av konsekvensane sine.

\section{Eit behov utan eit fag}

Skiljet skapte eit problem som handverket aldri hadde hatt: spørsmålet om kvaliteten på den industrielle forma. Så lenge snikkaren forma stolen i hand, var formkvaliteten innebygd i kompetansen hans; ein udugande snikkar laga få stolar lenge. Men når forma vart skild frå produksjonen og lagt i teikninga, vart formkvaliteten plutseleg eit ope spørsmål som ingen institusjon var ansvarleg for. Maskina ville lage kva som helst, like villig og like billeg, godt som dårleg. Teiknaren kunne teikne kva som helst. Og marknaden, som no var stor og namnlaus, kjøpte det som selde, ikkje det som var godt. For fyrste gong i historia kunne dårleg form produserast i enorme mengder og seljast med stor lønsemd, fordi ingen i kjeda korkje bar konsekvensen eller eigde standarden.

Dette er tomrommet faget seinare skulle bli kalla inn for å fylle. Nokon, ein eller annan stad, måtte ta ansvar for den industrielle forma, for at varene ikkje berre vart selde, men var verde å lage. Det er dette behovet, og ikkje ei indre fagleg modning, som kallar designfaget til verda. Men sjå tilstanden behovet oppstår i. Det finst eit reelt behov for nokon som kan svare for forma. Det finst ein ny skikkelse, teiknaren, som teknisk sett produserer forma. Men det finst ingen kunnskap som fortel korkje teiknaren eller nokon annan kva som gjer den industrielle forma god, fordi den einaste formkunnskapen som fanst, handverkets tause dømekraft, sat i hender som no var kopla ut av prosessen. Behovet er der. Skikkelsen er der. Fundamentet er ikkje der, og blei aldri grave ut; det vart tvert imot \emph{fylt att} i same rørsla, då tenkinga vart riven laus frå handa som åleine visste korleis ho skulle utøvast.

Det er rekkjefølgja som avgjer mykje av det som kjem. I dei modne faga går vegen den andre lei: fyrst veks det fram ein kunnskap, ein teori, ei rad observasjonar som heng saman, og så, av kunnskapen, eit fag som forvaltar han, og til slutt eit yrke som brukar han. Kjemikaren finst fordi kjemien finst; legen i moderne meining finst fordi fysiologien finst. Hjå design er rekkjefølgja snudd på hovudet. Fyrst kjem yrket, skapt av ein industri som trong nokon til å teikne. Så, langt seinare, kjem faget, kalla inn for å gje det alt eksisterande yrket eit grunnlag det aldri vart bygd på. Kunnskapen, fundamentet, skulle altså leverast \emph{etterskotsvis}, til ein praksis som alt var i full gang og alt definert av andre sine behov. Det er ein nær håplaus startposisjon, og han forklarer kvifor reformatorane som no kjem, og rørslene etter dei, alltid arbeider baklengs: dei prøver å destillere eit prinsipp ut av ein praksis som alt finst, i staden for å byggje ein praksis ut av eit prinsipp. Eit fag som må finne fundamentet sitt etter at huset er reist og bebudd, leitar sjeldan med same alvor som eit fag som måtte ha grunnen klar før det kunne byrje. Heile den vidare historia er soga om dette etterskotsvise søket, og om kor lett det er, når huset alt står, å pynte fasaden og kalle det å ha funne grunnen.

\section{Soho, og rommet der konseptet skilde seg frå handa}

Skiljet treng ikkje vere ein abstraksjon; det har golv, vegger og lys, og ein kan gå inn i rommet der det vart innreidd. Gå til Birmingham, til Matthew Boulton sin Soho Manufactory, opna i 1766: det store, palassliknande fabrikkbygget i Handsworth nord for byen, der opp mot sju hundre arbeidarar produserte knappar, spenner, lysestakar, urkasser og smågods i sølv og forsølva metall. Det vi skal sjå etter, er ikkje maskinene, men eit særskilt rom inne i bygget: teiknekontoret, der mønstera vart laga, og der ingen smelta metall eller stansa ein einaste knapp. Her sat folk og produserte former på papir, snitt, profilar, ornamentmønster, som så vart leverte ned til verkstadene der hender dei aldri rørte realiserte dei i tusental. Mellom desse to romma, det øvre der forma vart tenkt og det nedre der ho vart laga, gjekk skiljet vi har snakka om, no som ei fysisk trapp.

Boulton sjølv var ein kremmar av forma. Han henta inn mønster utanfrå, kopierte og tilpassa antikke motiv, og selde det same metallgodset i fleire prisklassar ved å variere finishen, ikkje sjølve forma. Då han slo seg saman med James Watt i 1775 og Soho seinare vart staden der den forbetra dampmaskina vart bygd, fekk anlegget eit nesten emblematisk tyngdepunkt i industrihistoria: her møttest krafta og forma, energien og pynten, under same tak. Men poenget vårt er ikkje storleiken eller dampen. Det er at Soho gjorde synleg, i si eiga romplan, det knappenålsfabrikken til Smith berre hadde skildra i ord. Smith fortalde om delinga; Boulton bygde henne i mur og innreidde henne med eit eige rom for dei som tenkte forma. John Heskett har peika på at det er i nett slike tidlege engelske manufakturar, metallvarene i Birmingham, keramikken i Staffordshire, at den utskilde formfunksjonen fyrst trer fram som ei eiga rolle i produksjonen, lenge før nokon kalla henne design \autocite{heskett1980}.

Sjå kva trappa gjorde med kunnskapen, for det er der det avgjerande skjer. Mannen i teiknekontoret kunne teikne ein lysestake som var elegant på papiret, men han kjende ikkje lenger frå si eiga hand korleis sølvet oppførte seg under stansa, kor tjukk kanten måtte vere for ikkje å skjere, kva som hende med eit tynt parti når det vart polert hundre gonger. Den kunnskapen sat no nede i verkstaden, i hendene til dei som faktisk laga tingen, men dei tok ikkje formavgjerdene. Og mannen som tok formavgjerdene, hadde ikkje kunnskapen i handa lenger. Adrian Forty har formulert dette presist: forma vart skild frå produksjonen ikkje av tekniske grunnar, men fordi det gav eigaren kontroll over ein kunnskap han elles var prisgjeven hjå dyre handverkarar \autocite{forty1986}. På Soho ser ein prisen for den kontrollen reint. Konseptet steig opp ei trapp, og handa vart att nede. Og frå den dagen var det to slags menneske i bygget: dei som visste korleis forma skulle sjå ut utan å kunne lage henne, og dei som kunne lage henne utan å få bestemme korleis ho skulle sjå ut. Verken den eine eller den andre eigde forma heil.

Det uskuldige draget er det som gjer skiljet varig. Ingen på Soho opplevde at dei reiv noko frå kvarandre. Dei rasjonaliserte ein produksjon, dei laga vakrare og billegare gods, dei opna ein eksport. Teiknekontoret var ein effektivitet, ikkje eit overgrep. Men sett frå kunnskapssida var det ei amputasjon utført så reint at pasienten ikkje merka det: den tause dømekrafta som hadde butt i éin kropp, vart delt i to halvdelar som aldri meir skulle møtast i éin person. Boulton kunne no eige forma som ei vare på lager med mønster, uavhengig av kven som sto ved stansa den dagen. Det var heile sigeren, og det er denne sigeren, ikkje noko nederlag, designfaget seinare arvar: ein formfunksjon som frå fyrste stund var rissa inn i ei trapp mellom to rom, der den eine enden hadde auget og den andre handa, og ingen hadde begge.

\section{Birmingham-knappen, og forma som mista heile følgjet sitt}

Følg éin einaste gjenstand heile vegen, for det abstrakte skiljet mellom form og følgje får fyrst tyngd når ein går strekninga sjølv, ledd for ledd. Vel ein metallknapp, laga i Birmingham ein gong i 1790-åra. I handverkets verd hadde ein knapp eit kort liv mellom kjende hender: ein lokal handverkar laga han, ein lokal kunde bar han, og om han svikta, om plateringa flassa, om festet brast, fekk handverkaren det att frå mannen han møtte på torget. Forma, lagingane og følgja budde innanfor ein dags gange. Den industrielle knappen lever eit heilt anna liv. Han vert teikna i eit kontor av ein som aldri ser kven som ber han; stansa ut i tusental av ei hand som ikkje teikna han; forsølva i eit bad av ein arbeidar som pustar dampen frå syrene; pakka, send, seld i ein by teiknaren aldri har vore i, til ein kjøpar han aldri møter; og når plateringa flassar etter eit halvt år, fell den følgja på ein framand, langt utanfor synsranda til kvar einaste hand som var med å lage han.

Tel ledda. Mellom den som bestemte forma og den som bar henne ligg det ein teiknar, ein stansar, ein platerar, ein pakkar, ein grossist, ein kjøpmann og ein butikkmann, eit halvt dusin menneske og fleire hundre kilometer. At plateringa er for tynn til å vare, melder seg hjå kjøparen om eit halvt år. At syrebadet øydelegg lungene, melder seg hjå plateraren om ti år. At metallet ikkje kan brukast om att, melder seg på avfallsdyngja. Ingen av desse meldingane når attende til teiknebordet, fordi strukturen er bygd slik at dei ikkje skal nå det. Teiknaren i Birmingham er ikkje ein dårleg mann; han er ein mann plassert utanfor rekkjevidde av konsekvensane sine av sjølve arbeidsdelinga han står i, og han kan ikkje bere ei følgje han er strukturelt skild frå. Det han stod oppi, var ei romplan før det var ei skuld: teiknekontoret over verkstaden plasserte den som forma forma utanfor staden der forma fall, og rommet tok avgjerda for han lenge før han sjølv kunne ta henne. David Pye sitt skilje får her sin sosiale kant: visseheitens handverk fjerna ikkje berre risikoen frå utføringa ved hjelp av malen og maskina, det fjerna òg følgja frå synsranda til den som teikna malen.

\section{Maskinbrytarane, og motstanden mot skiljet}

Det er lett å lese skiljet som ein lydlaus prosess, ei administrativ omlegging utan offer som protesterte. Det var det ikkje. Skiljet vart òg slege, med slegge, om natta, i lyset frå brennande verkstader, og dei som slo visste betre enn dei fleste seinare teoretikarar kva som var i ferd med å bli teke frå dei. Mellom mars 1811 og 1816 gjekk det ei bylgje av maskinøydeleggingar over Nottinghamshire, Yorkshire og Lancashire, strømpevevarane og klypparane som ettertida kalla luddittane, etter den halvt oppdikta «General Ludd» dei skreiv brev i namnet til. Dei braut seg inn i verkstader om natta og slo sund dei breie vevstolane og klyppemaskinane, ikkje vilt, men utvalt: dei lét gamle maskinar stå og knuste dei nye, dei som let ein ufaglærd gut gjere på timar det ein opplærd handverkar hadde brukt år på å meistre.

Det var ikkje maskina som sådan dei slo etter. Strømpevevaren i Nottingham hadde sjølv vore ein mann med ei maskin i stova; han hata ikkje verktøyet. Det han slo etter, var den nye breie ramma som gjorde at forma, mønsteret, snittet, kvaliteten på maska, vart fastsett ein annan stad, av ein annan, og at han sjølv vart redusert til ein som mata stoff inn i ei avgjerd han ikkje lenger tok del i. Han slo etter sjølve skiljet. Han hadde vore ein som tenkte forma og laga henne i same rørsle, og no skulle han bli ein som berre laga, til lægre løn, av dårlegare garn, etter eit mønster nokon i eit kontor hadde bestemt. Bylgja vart slått ned med ein hardleik som seier kor mykje som stod på spel: i februar 1812 gjorde Parlamentet maskinbryting til ei forbrytelse som kunne straffast med døden, Byron heldt si fyrste tale i Overhuset mot lova, og i januar 1813 vart sytten menn hengde i York etter ein einaste rettargang. Staten verna det nye skiljet med galgen.

Den seinare designhistoria har ikkje hatt rom for desse mennene, og det er sigande. Dei passar ikkje inn, fordi dei stod på feil side av det skiljet faget seinare skulle gjere til sjølve grunnvollen sin, sida til handa, ikkje sida til teikninga. Frank Trentmann har minna om at framvoksteren av varesamfunnet ikkje var ein roleg straum, men ein strid med taparar som visste at dei tapte \autocite{trentmann2016}. Luddittane er det fyrste klare dømet på nokon som såg skiljet medan det vart innført, og kalla det ved namn: ikkje framsteg, men ekspropriasjon av kunnskapen frå handa som åtte han. At dei tapte, at det breie vevstolen vann, at forma flytta til kontoret, at galgen fekk siste ordet, er ikkje eit prov på at dei tok feil om kva som hende. Det er berre eit prov på kven som hadde makta. Og det legg ein mørk understrek under heile boka sitt poeng: faget vart fødd på sigerssida i ein strid det aldri sidan har villa hugse at vart utkjempa.

\section{Babbage, og delinga ført i rekneskap}

Tjue år etter at galgane stod i York, gav ein matematikar i London skiljet det språket det hadde mangla, eit reint, kaldt, økonomisk språk som la heile logikken open. I 1832 kom Charles Babbage, mannen som elles er hugsa for dei mekaniske reknemaskinane sine, ut med \textit{On the Economy of Machinery and Manufactures}, ei bok som ikkje handla om form i det heile, men som forklarar den industrielle forma sin lagnad betre enn nokon estetisk avhandling i hundreåret. Babbage gjekk gjennom fabrikk etter fabrikk med eit reknebrett i handa og spurde, om kvar einaste operasjon: kor mykje dyktigheit krev ho, og kor mykje må eg betale for den dyktigheita? Av desse spørsmåla destillerte han det som ettertida kallar «Babbage-prinsippet», og det er prinsippet under heile skiljet, no skrive som rekneskap.

Sjå nøye på kva det seier, for det er nådelaust klårt. Når ein deler ein arbeidsprosess i åtskilde operasjonar, treng arbeidsgjevaren berre kjøpe akkurat så mykje dyktigheit kvar operasjon krev, og ikkje meir. I det udelte handverket måtte ein betale den fullt opplærde mannen for heile dagen, også for dei timane han gjorde enkelt arbeid ein gut kunne gjort, fordi han var éin kropp og kunne ikkje delast. Deler ein arbeidet, kan ein setje guten til det enkle og spare den dyre mannen til det vanskelege. Babbage rekna det ut på ei knappenål, akkurat som Smith, men der Smith såg produktiviteten, såg Babbage prisen på kunnskap. Forma, det vanskelegaste, det som krev mest dyktigheit, det som ein gong budde i den fullt opplærde handa, blir i denne reknemåten den dyraste operasjonen, og difor den ein for all del må skilje ut og kjøpe åtskilt, i minst mogeleg mengd, frå éin person, slik at ein slepp å betale for henne i kvar einaste arbeidar. John Heskett har peika på at det er denne økonomiske logikken, ikkje noko estetisk program, som driv den utskilde formfunksjonen fram i den tidlege industrien \autocite{heskett1980}.

Her ligg det avgjerande for boka, og Babbage seier det utan å vite at han seier det. Når delinga er ein rekneskap, blir forma ein \emph{kostnad} før ho er noko anna. Ho er den dyre operasjonen, den ein kjøper minst av og skil lengst ut. Designaren, teiknaren, modelløren, er i Babbage sin reknemåte rett og slett den dyraste arbeidaren, han som ber den dyktigheita ein ikkje vil betale for i mengd, og som difor må sentraliserast i éin og leverast vidare til hender som er billegare fordi dei berre utfører. Adrian Forty har lese den same logikken i fabrikkrekneskapen frå perioden: forma vart skild frå produksjonen fordi det stod i bøkene at det lønte seg \autocite{forty1986}. Og legg merke til kva dette gjer med spørsmålet om eit fundament. Eit fag som frå fyrste rekneskapslinje er definert som ein kostnad å minimere, blir aldri spurt om kva forma \emph{burde} vere, berre om kor lite av denne dyre evna ein kan klare seg med. Penny Sparke har understreka kor djupt den økonomiske rasjonaliteten er innvoven i sjølve framvoksteren av designyrket \autocite{sparke2004}. Babbage gav skiljet rekneskapen sin, og rekneskapen hadde inga linje for kva forma var verd, berre for kva ho kosta.

\section{Den som teiknar og den som ber følgja}

Det er ein siste konsekvens av delinga som dei to skikkelsane, luddittane og Babbage, saman gjer synleg. Delinga gjorde ikkje berre forma til ein kostnad og kunnskapen til noko ein person eigde og ein annan utførte. Ho gjorde forma til noko \emph{éin person teikna og ein heilt annan bar konsekvensen av}, og dei to møttest aldri. I det udelte handverket fall følgja av ei dårleg avgjerd attende på den som tok henne: laga snikkaren ein vond stol, var det han som mista kunden og han som måtte gjere arbeidet om. Forma og følgja budde i same kropp, og difor heldt følgja forma i tøyle.

Delinga reiv dei frå kvarandre, ledd for ledd, og kvart ledd flytta følgja eit steg lenger frå den som teikna. Mannen i kontoret teikna det breie mønsteret; strømpevevaren bar tapet av yrket sitt. Babbage teikna rekneskapen; guten ved maskina bar dei lange timane og den låge løna han hadde rekna ut. Mønsterteiknaren la profilen; arbeidaren i syrebadet bar lungene. Og kjøparen, langt nede i kjeda, bar den dårlege maska, den tynne plateringa, den korte levetida som teiknaren aldri såg fordi han var skild frå staden der ho synte seg. Dette er ikkje fire ulike historier. Det er éi: delinga plasserte avgjerda om forma i éin person og følgja av henne i mange andre, og bygde avstanden mellom dei så stor at følgja ikkje lenger kunne nå attende og rette avgjerda. Frank Trentmann har vist kor systematisk denne avstanden mellom avgjerd og følgje voks gjennom heile den industrielle epoken \autocite{trentmann2016}. Han voks ikkje ved eit uhell. Han voks fordi han lønte seg, og fordi ingen i kjeda hadde både handa som teikna og ryggen som bar.

\vignett

Inn i dette tomrommet skulle det snart stige menn som meinte dei visste svaret. Dei såg den industrielle forma, billeg, overpynta, laga for sal, og dei fann henne stygg og umoralsk, og dei sette seg føre å reformere henne. Dei kom med prinsipp, med katalogar, med ein moral om ornament og ærlegdom. For fyrste gong prøvde nokon å gje den industrielle formgjevinga eit grunnlag. Det fyrste forsøket har ein bygning og eit år, og det er dit vi går no.

\laeringsmal{\begin{itemize}
\item Forklare korleis arbeidsdelinga skil forma frå produksjonen, og kvifor dette skiljet var sosialt og økonomisk før det vart teknisk.
\item Gjere greie for skilnaden mellom skiljet form/produksjon og det alvorlegare skiljet form/følgje, og kvifor det siste plasserer formgjevaren utanfor rekkjevidde av konsekvensane sine.
\item Bruke Wedgwood og Etruria som konkret døme på korleis «designaren» oppstod som funksjonær, ikkje som fagperson.
\item Grunngje kvifor design fekk eit yrke før det fekk eit fag, og kva denne snudde rekkjefølgja har å seie for alt som kjem etter.
\end{itemize}}

\ovingar{\begin{enumerate}
\item Wedgwood ville «make such machines of the men as cannot err». Diskuter setninga opp mot Pye sitt skilje mellom risikoens og visseheitens handverk. Kva gjekk tapt, og kva vart vunne?
\item Vel ein gjenstand du eig. Prøv å spore kjeda frå den som bestemte forma til den som bar konsekvensen av henne. Kor mange ledd ligg mellom, og kva gjer avstanden med ansvaret?
\item «Eit yrke utan eit fag er nett det skiljet etterlét.» Samanlikn design med eit anna yrke som voks fram av industrien (til dømes maskinist eller fotograf). Fekk dei eit fag? Korleis?
\item Boka hevdar at fundamentet skulle leverast «etterskotsvis» til ein praksis som alt var i gang. Kan eit fag byggjast baklengs på denne måten? Drøft med eit døme der det lukkast og eit der det ikkje gjorde.
\end{enumerate}}

\vidarelesing{\fullcite{forty1986}; \fullcite{heskett1980}; \fullcite{sparke2004}; \fullcite{trentmann2016}; \fullcite{pugin1836contrasts}; \fullcite{ruskin1853stones}; \fullcite{banham1960}.}
